% -*- coding: utf-8 -*-
% !TEX program = xelatex
\documentclass[plain,noanswer,medmath]{../../template/jnuexam}
\usepackage{../../template/kysx}

\begin{document}


\examtitle{1988}{5}


\loadproblems{3}{1988P3}%载入试卷三
\loadproblems{4}{1988P4}%载入试卷四


\makepart{填空题}{本题满分12分，每空1分}%【同试卷四第一题】


\useproblem{4}{1}{1}


\useproblem{4}{1}{2}


\useproblem{4}{1}{3}


\useproblem{4}{1}{4}


\makepart{判断题}{本题满分10分，每小题2分}%【同试卷四第二题】
%本题满分10分；每小题回答正确得2分，回答错误得-1分，不回答得0分；全题最低得0分


\useproblem{4}{2}{1}


\useproblem{4}{2}{2}


\useproblem{4}{2}{3}


\useproblem{4}{2}{4}


\useproblem{4}{2}{5}


\makepart{计算题}{本题满分16分，每小题4分}


\begin{problem}
求极限$\lim_{x \to 1} (1 - x^2)\tan \frac{\pi x}{2}$．
\end{problem}


\begin{problem}
已知$u = \e^{\frac{x}{y}}$，求$\frac{\pd^2 u}{\pdx\pd y}$．
\end{problem}


\useproblem{4}{3}{3}%【同试卷四第三(3)题】


\useproblem{4}{3}{4}%【同试卷四第三(4)题】


\makepart{}{本题满分6分}

\begin{problem}
确定常数$a$和$b$，使函数$f(x) = \begin{cases}
  ax + b, & x > 1, \\
     x^2, & x \le 1,
\end{cases}$处处可导．
\end{problem}


\makepart{}{本题满分8分}%【同试卷三第五题】

\useproblem{3}{5}{0}%


\makepart{}{本题满分8分}%【同试卷四第六题】

\useproblem{4}{6}{0}%


\makepart{}{本题满分8分}%【同试卷四第七题】

\useproblem{4}{7}{0}%


\makepart{}{本题满分6分}

\begin{problem}
已知$n$阶方阵$A$满足矩阵方程$A^2 - 3A - 2E = O$，其中$A$是给定的，而$E$是单位矩阵．
证明$A$可逆，并求了其逆矩阵$A^{-1}$．
\end{problem}


\makepart{}{本题满分7分}%【同试卷四第八题】

\useproblem{4}{8}{0}%


\makepart{}{本题满分7分}%【同试卷四第十题】

\useproblem{4}{10}{0}


\makepart{}{本题满分7分}

\begin{problem}
假设有十只同种电器元件，其中有两只废品，装配仪器时，从这批元件中任取一只，
如是废品，则扔掉重新任取一只；如仍是废品，则扔掉再取一只．
试求在取到正品之前，已取出的废品只数的分布、数学期望和方差．
\end{problem}


\makepart{}{本题满分5分}%【同试卷四第十二题】

\useproblem{4}{12}{0}


\end{document}
